\section{Discussion}
\label{sec:discussion}

\paragraph{When the defense matters most.}
Theorem~\ref{thm:per-client} decomposes per-client leakage into a
controllable mechanism term and an uncontrollable prior-coupling floor.
The defense can only shrink the controllable term; the floor is
deployment-determined and beyond DP-SGD's reach. The defense's
\emph{relative} value is therefore highest when the controllable term
dominates — \ie when DP noise is calibrated to be informative
($\sigma$ small enough that $a/\sigma_i^2$ is comparable to
$\ell_i^\circ$). Once $\ell_i^\circ$ swamps $a/\sigma_i^2$, additional
noise allocation produces diminishing returns; the appropriate response
is then either to reduce the prior coupling (change the partitioning,
remove organisational metadata from the threat model) or to introduce
secure aggregation~\cite{bonawitz2017secagg} as a complementary
defense.

\paragraph{The bound is conservative in a deployment-favourable
direction.}
Theorem~\ref{thm:per-client}'s prior-coupling term $\ell_i^\circ$ is
defined as a \emph{supremum} over the family of priors
$\FAM_{\TOPO,\omega}$: it is the worst case across all priors
consistent with the deployment's topology and organisational
structure. Achieving the supremum empirically requires the
adversary's shadow training distribution to match the target's
deployment prior. \S\ref{sec:experiments-channels} provides direct
empirical evidence on this point. In Setting~C, where shadow and
target use the same $\eta$-coupling adapter (same prior family),
$\ell_i^\circ$ is empirically realised: the org channel attack lift
is positive, scales monotonically with $\eta$, and vanishes at
$\eta = 0$. In Settings~B and~A, where the adversary has only
synthetic Dirichlet re-partitioning of public data while the target
uses native FLamby cross-silo sites, three orthogonal
methodological improvements (DP-matched shadow, class-aware
features, cross-validated regressor) moved $\Adv_2^{\text{full}}$
AUROC across the 0.5 threshold (0.19~$\to$~0.56) but did not
flip attack lift positive. The supremum is not realised because the
adversary's prior does not match.

This is the right kind of conservatism for a privacy guarantee.
Fulcrum's bound is valid as an upper bound across all priors in the
family; it is tight when the adversary's prior matches the
deployment's; it is loose --- in a \emph{deployment-favourable}
direction --- when the adversary has only public-proxy data without
site labels. The bound therefore over-states realisable attack
capability in cross-silo deployments where adversaries cannot
obtain matching shadow priors, which is the typical condition for
hospitals, banks, and industrial consortia operating under
data-sharing constraints. We do not weaken the bound to remove this
slack: privacy guarantees are appropriately conservative when the
threat-model assumption (the worst-case adversary in the prior
family) is stronger than the realistic adversary (a public-proxy
attacker with no site-labelled training data). The empirical
characterisation in \S\ref{sec:experiments-channels} \emph{quantifies}
this conservatism rather than claiming the bound is tight.

\paragraph{Why our defense reduces to uniform under symmetry.}
Corollary~\ref{cor:strict} establishes that the topology-aware allocation
collapses to standard uniform DP-SGD whenever leverage scores are
identical across clients. We treat this as a feature: the theorem
recovers the standard practice exactly in the regime that the standard
practice was designed for, and then strictly improves on it whenever
the deployment introduces structural asymmetry. This yields a clean
deployment recommendation: \emph{use \Fulcrum unconditionally} — when
the topology happens to be symmetric, you pay no price for the choice.

\paragraph{Three proxies, three deployments.}
The three leverage proxies of \S\ref{sec:defense-proxies} are not
competing instantiations; they are appropriate for different regimes
that share a common structural language. Hierarchical FL with uneven
groups uses Corollary~\ref{cor:sbm}; decentralised FL with non-symmetric
topology uses Corollary~\ref{cor:degree}; cross-silo FL with
heterogeneous client dataset sizes uses
Corollary~\ref{cor:influence}. A practical deployment can blend
proxies by passing weighted leverage vectors to
Theorem~\ref{thm:allocation}. We do not address the question of
\emph{which} proxy is universally best because there is no universal
answer: the right proxy is the one that captures the dominant
asymmetry in the deployment's structural prior.

\paragraph{Honest limitations.}
Three limitations bound the scope of our claims.

\textit{Conservative bounds.}
Theorem~\ref{thm:per-client} drops a non-negative subtractive term in
the chain-rule decomposition (Appendix~\ref{app:thm1}) and uses the
Cuff-Yu KL-to-MI conversion~\cite{cuff2016mi} which is loose at extreme
$\sigma$. The bound is a valid upper bound but not tight; tighter
conversion via Asoodeh et~al.~\cite{asoodeh2021variants} or amplified
RDP via Wang et~al.~\cite{wang2019subsampled} would shrink the bound
substantially under Poisson subsampling. The conservative bound is
appropriate for a privacy guarantee (we want an upper bound on adversary
leakage, not the actual achievable leakage). Tighter bounds are a
follow-up direction.

\textit{Heuristic proxies.}
Corollaries~\ref{cor:degree} and~\ref{cor:influence} are heuristic; only
Corollary~\ref{cor:sbm} carries a rigorous information-theoretic
derivation. The heuristic proxies are validated empirically on real and
synthetic deployments, but a tight analytical bound for either remains
an open problem.

\textit{Disjoint datasets.}
Assumption (IA2) precludes deployments where the same individual
contributes data to multiple clients. Cross-device FL with overlapping
contributors (\eg a single user logged into multiple devices) violates
this assumption; the bound becomes loose in proportion to the overlap.
Cross-silo healthcare and industry consortia, the deployments we
target, naturally satisfy (IA2) because each site enrols disjoint
populations.

\paragraph{Cross-pollination with active threats.}
Our threat model is passive. Active adversaries that manipulate the
training protocol — \eg malicious aggregators that selectively retain
gradients, or Sybil clients that inject misleading updates — operate in
a regime that DP-SGD fundamentally cannot bound. We expect the
topology-aware allocation to be \emph{compatible} with active-defense
mechanisms (Byzantine-resilient aggregators, anomaly detection) in the
same way that DP-SGD is, but proving non-trivial guarantees in the
active setting requires distinct machinery.

\paragraph{Toward tighter theory.}
The natural next theoretical step is to prove the Wang/Balle/Kasiviswanathan
amplified RDP~\cite{wang2019subsampled} version of
Theorem~\ref{thm:per-client}. Under Poisson subsampling at rate
\(q = |B|/|\Dset_i|\), the per-round RDP scales as $q^2 \alpha/(2\sigma_i^2)$
rather than $\alpha/(2\sigma_i^2)$. The corresponding bound would be
\(\Tmax q_i^2 C^2 / (2\sigma_i^2 |B|^2) + \ell_i^\circ\), giving
\emph{client-specific} amplification: clients with smaller datasets
(higher $q_i$) gain more from amplification. This integrates naturally
with Corollary~\ref{cor:influence}'s dataset-size leverage proxy and
deserves a dedicated treatment.
